% In this file you should put the actual content of the blueprint.
% It will be used both by the web and the print version.
% It should *not* include the \begin{document}
%
% If you want to split the blueprint content into several files then
% the current file can be a simple sequence of \input. Otherwise It
% can start with a \section or \chapter for instance.

% content.tex  —  leanblueprint for Pauli Weight Truncation Error Bounds
% Based on: "Bounds Beating: Pauli Weight Truncation Errors" (Teng, after Angrisani et al. 2025)

\chapter{Pauli Weight Truncation: Setup and Definitions}

\section{Pauli Operators and Weight}

\begin{definition}[Pauli group and weight]
  \label{def:pauli-group}
  \lean{PauliGroup}
  \leanok
  The \emph{$n$-qubit Pauli group} $\mathcal{P}_n$ consists of all $n$-fold tensor
  products of single-qubit Paulis $\{I, X, Y, Z\}$.  For $P \in \mathcal{P}_n$, the
  \emph{weight} $|P|$ is the number of tensor factors that are not the identity $I$.
\end{definition}

\begin{definition}[Normalized Pauli basis]
  \label{def:normalized-pauli-basis}
  \lean{normalizedPauliBasis}
  \leanok
  The \emph{normalized Pauli basis} is
  \[
    \widetilde{\mathcal{P}}_n
    \;=\;
    \Bigl\{\tfrac{I}{\sqrt{2}},\, \tfrac{X}{\sqrt{2}},\, \tfrac{Y}{\sqrt{2}},\,
            \tfrac{Z}{\sqrt{2}}\Bigr\}^{\!\otimes n}.
  \]
  Its elements satisfy $\operatorname{Tr}[s\,t] = \delta_{st}$ for all
  $s, t \in \widetilde{\mathcal{P}}_n$.
\end{definition}

\begin{definition}[Normalized Pauli norm]
  \label{def:pauli-norm}
  \lean{pauliNorm}
  \uses{def:normalized-pauli-basis}
  \leanok
  For an observable $O = \sum_P a_P P$ expanded in the Pauli basis, its
  \emph{normalized Hilbert--Schmidt (Pauli) norm} is
  \[
    \|O\|_{\mathrm{Pauli},2}
    \;:=\;
    \Bigl(\sum_P a_P^2\Bigr)^{1/2}
    \;=\; 2^{-n/2}\,\|O\|_2.
  \]
\end{definition}

\section{Locally Scrambling Distributions}

\begin{definition}[Locally scrambling layer]
  \label{def:locally-scrambling}
  \lean{LocallyScrambling}
  \leanok
  A random unitary layer $U_j$ is \emph{locally scrambling} if its distribution is
  invariant under conjugation by independent random single-qubit Clifford gates on
  each qubit.
\end{definition}

\begin{lemma}[Pauli orthogonality]
  \label{lem:pauli-orthogonality}
  \lean{pauliOrthogonality}
  \uses{def:locally-scrambling, def:pauli-group}
  \leanok
  Let $U_j$ be locally scrambling. Then for any $P \neq Q$ in $\mathcal{P}_n$,
  \[
    \mathbb{E}_{U_j}\!\left[U_j^{\dagger\otimes 2}(P \otimes Q)\,U_j^{\otimes 2}\right] = 0.
  \]
\end{lemma}

\begin{lemma}[Pauli mixing inequality]
  \label{lem:pauli-mixing}
  \lean{pauliMixing}
  \uses{def:locally-scrambling, def:pauli-group}
  \leanok
  Let $U_j$ be locally scrambling and let $\rho$ be any $n$-qubit state.  For any
  non-identity $P \in \mathcal{P}_n$,
  \[
    \mathbb{E}_{U_j}\!\left[\operatorname{Tr}\!\bigl(P\,U_j\rho U_j^\dagger\bigr)^2\right]
    \;\leq\;
    \Bigl(\tfrac{2}{3}\Bigr)^{|P|}.
  \]
\end{lemma}

\begin{definition}[Shallow layer assumption]
  \label{def:shallow-layer}
  \lean{ShallowLayer}
  \uses{def:locally-scrambling, def:pauli-group}
  Each layer $U_j$ is \emph{shallow} if, for every $P$ of weight $|P| = k$, the
  Heisenberg-evolved operator $U_j^\dagger P\, U_j$ expands into at most
  $n^{\mathcal{O}(k)}$ Pauli terms.
\end{definition}

\chapter{The Weight Truncation Algorithm}

\section{Recursive Definition of Truncated Observables}

\begin{definition}[Truncated observable sequence]
  \label{def:truncated-observable}
  \lean{truncatedObservable}
  \uses{def:normalized-pauli-basis, def:pauli-group}
  \leanok
  Let $U = U_L \cdots U_1$ be an $L$-layer circuit, $O$ an observable, and $k \geq 0$
  an integer threshold.  Define a sequence $O_{L+1} := O$ and, recursively for
  $j = L, L-1, \ldots, 1$,
  \[
    O_j
    \;:=\;
    \sum_{\substack{s \in \widetilde{\mathcal{P}}_n \\ |s| \leq k}}
    \operatorname{Tr}\!\bigl[U_{j+1}^\dagger O_{j+1} U_{j+1} \cdot s\bigr]\,s.
  \]
  The \emph{weight-$k$ truncated observable} is $O_U^{(k)} := O_1$.
\end{definition}

\begin{definition}[Truncated estimator]
  \label{def:truncated-estimator}
  \lean{truncatedEstimator}
  \uses{def:truncated-observable}
  \leanok
  Given an input state $\rho$ and a circuit $U$, the \emph{truncated estimator} is
  \[
    \tilde{f}_U^{(k)}(O) \;:=\; \operatorname{Tr}\!\bigl[O_U^{(k)}\,\rho\bigr].
  \]
  The true expectation value is $f_U(O) := \operatorname{Tr}[O\,U\rho U^\dagger]$.
\end{definition}

\section{Telescoping Decomposition of the Error}

\begin{lemma}[Telescoping error decomposition]
  \label{lem:telescoping}
  \lean{telescopingError}
  \uses{def:truncated-estimator, def:truncated-observable}
  The mean squared error admits a layer-by-layer decomposition
  \[
    \mathbb{E}_U\!\bigl[\bigl(f_U(O) - \tilde{f}_U^{(k)}(O)\bigr)^2\bigr]
    \;=\;
    \sum_{j=1}^{L}
    \mathbb{E}_U\!\Bigl[
      \operatorname{Tr}[O_{j+1}\rho_{j+1}]^2
      - \operatorname{Tr}[O_j \rho_j]^2
    \Bigr],
  \]
  where $\rho_j = U_j \cdots U_1\,\rho\,U_1^\dagger \cdots U_j^\dagger$.
\end{lemma}

\begin{lemma}[Cross-term vanishing]
  \label{lem:cross-terms}
  \lean{crossTermsVanish}
  \uses{lem:pauli-orthogonality, lem:telescoping}
  Under the locally scrambling assumption, the cross-terms in the second-moment
  expansion of each summand in Lemma~\ref{lem:telescoping} vanish, i.e.\
  \[
    \mathbb{E}_{U_j}\!\left[
      \operatorname{Tr}[s\,U_j\rho_j U_j^\dagger]\,
      \operatorname{Tr}[t\,U_j\rho_j U_j^\dagger]
    \right] = 0
    \qquad \text{for } s \neq t \in \widetilde{\mathcal{P}}_n.
  \]
\end{lemma}

\begin{lemma}[Per-layer error bound]
  \label{lem:per-layer}
  \lean{perLayerBound}
  \uses{lem:pauli-mixing, lem:cross-terms, def:pauli-norm}
  Each non-negative summand in the telescoping decomposition satisfies
  \[
    \mathbb{E}_U\!\Bigl[
      \operatorname{Tr}[O_{j+1}\rho_{j+1}]^2 - \operatorname{Tr}[O_j\rho_j]^2
    \Bigr]
    \;\leq\;
    \Bigl(\tfrac{2}{3}\Bigr)^{k+1}\,\|O\|_{\mathrm{Pauli},2}^2.
  \]
\end{lemma}

\chapter{Main Result}

\begin{theorem}[Mean squared truncation error — depth-independent bound]
  \label{thm:mse-bound}
  \lean{mseBound}
  \uses{lem:telescoping, lem:cross-terms, lem:per-layer,
        def:locally-scrambling, def:truncated-estimator, def:pauli-norm}
  % \leanok
  Let $U = U_L \cdots U_1$ be drawn from an $L$-layered locally scrambling circuit
  distribution.  For any integer $k \geq 0$, any observable $O$, and any $n$-qubit
  state $\rho$,
  \[
    \mathbb{E}_U\!\left[\bigl(f_U(O) - \tilde{f}_U^{(k)}(O)\bigr)^2\right]
    \;\leq\;
    \Bigl(\tfrac{2}{3}\Bigr)^{k+1}\,\|O\|_{\mathrm{Pauli},2}^2.
  \]
  In particular, the bound is independent of the circuit depth $L$.
\end{theorem}

\begin{proof}
  Apply Lemma~\ref{lem:telescoping} to write the MSE as a sum over layers.
  Lemma~\ref{lem:cross-terms} (using local scrambling via
  Lemma~\ref{lem:pauli-orthogonality}) kills all cross-terms in the second-moment
  expansion.  Lemma~\ref{lem:pauli-mixing} bounds each diagonal term by
  $(2/3)^{k+1}$ times the relevant squared Pauli coefficient, and summing over
  layers yields a geometric series whose value is again $(2/3)^{k+1}$ (the
  depth-$L$ prefactor cancels).  Combining with the definition of
  $\|\cdot\|_{\mathrm{Pauli},2}$ gives the result.
\end{proof}

\chapter{Conjectured Depth-Dependent Improvement}

\begin{conjecture}[Depth-decaying MSE bound]
  \label{conj:depth-decay}
  \lean{depthDecayConjecture}
  \uses{thm:mse-bound, def:locally-scrambling}
  There exists a universal constant $c \in (0,1)$ such that, under the same
  hypotheses as Theorem~\ref{thm:mse-bound},
  \[
    \mathbb{E}_U\!\left[\bigl(f_U(O) - \tilde{f}_U^{(k)}(O)\bigr)^2\right]
    \;\leq\;
    \Bigl(\tfrac{2}{3}\Bigr)^{k+1}\,c^L\,\|O\|_{\mathrm{Pauli},2}^2.
  \]
\end{conjecture}

\begin{remark}
  Conjecture~\ref{conj:depth-decay} would imply that deeper circuits are easier to
  classically shadow: the truncated estimator becomes more accurate as $L$ grows,
  at an exponential rate in the depth.  Proving this would require tracking
  cancellations between layers rather than bounding each layer independently.
\end{remark}
